% Part: many-valued-logics
% Chapter: sequent-calculus
% Section: proof-theoretic-notions

\documentclass[../../../include/open-logic-section]{subfiles}

\begin{document}

\olfileid[id]{mvl}{seq}{prf}

\olsection{Gagasan Teori-Bukti}

\begin{explain}
Sebagaimana kita telah mendefinisikan sejumlah gagasan semantik penting
(validitas, perikutan, keterpenuhan), kini kita mendefinisikan
\emph{gagasan teori-bukti} yang bersesuaian. Gagasan-gagasan ini tidak
didefinisikan dengan mengacu pada terpenuhinya !!{sentence} di dalam
!!{structure}, melainkan dengan mengacu pada !!{derivability} atau
!!{nonderivability} sekuen-sekuen tertentu. Penemuan bahwa gagasan-gagasan
tersebut berimpit merupakan penemuan penting. Keberimpitan itu adalah isi
\emph{teorema kesahihan} dan \emph{teorema kelengkapan}.
\end{explain}

\begin{defn}[Teorema]
!!^a{sentence}~$!A$ adalah sebuah \emph{teorema} jika terdapat
!!a{derivation} dalam~$\Log{LK}$ untuk sekuen $\quad \Sequent !A$.
Kita menulis $\Proves !A$ jika $!A$ merupakan teorema dan
$\Proves/ !A$ jika bukan.
\end{defn}

\begin{defn}[!!^{derivability}]
!!^a{sentence}~$!A$ \emph{!!{derivable} dari} suatu himpunan
!!{sentence}~$\Gamma$, $\Gamma \Proves !A$, jika dan hanya jika terdapat
himpunan bagian berhingga~$\Gamma_0 \subseteq \Gamma$ dan suatu barisan
$\Gamma_0'$ dari !!{sentence} di~$\Gamma_0$ sedemikian sehingga
$\Log{LK}$ !!{derive} $\Gamma_0' \Sequent !A$. Jika $!A$ tidak
!!{derivable} dari $\Gamma$, kita menulis $\Gamma \Proves/ !A$.
\end{defn}

Karena adanya aturan kontraksi, pelemahan, dan pertukaran, urutan dan
banyaknya !!{sentence} di dalam~$\Gamma_0'$ tidaklah penting: jika sekuen
$\Gamma_0' \Sequent !A$ !!{derivable}, maka begitu pula
$\Gamma_0'' \Sequent !A$ untuk setiap $\Gamma_0''$ yang memuat
!!{sentence} yang sama dengan~$\Gamma_0'$. Sebagai contoh, jika
$\Gamma_0 = \{!B, !C\}$, maka baik $\Gamma_0' = \tuple{!B, !B, !C}$
maupun $\Gamma_0'' = \tuple{!C, !C, !B}$ merupakan barisan yang hanya
memuat !!{sentence} di~$\Gamma_0$. Jika sekuen yang memuat salah satu
barisan itu !!{derivable}, maka demikian pula sekuen yang memuat barisan
lainnya, misalnya:
\begin{prooftree}
  \AxiomC{}
  \Deduce$!B, !B, !C \fCenter !A$
  \RightLabel{\LeftR{\Contraction}}
  \UnaryInf$!B, !C \fCenter !A$
  \RightLabel{\LeftR{\Exchange}}
  \UnaryInf$!C, !B \fCenter !A$
  \RightLabel{\LeftR{\Weakening}}
  \UnaryInf$!C, !C, !B \fCenter !A$
\end{prooftree}
Mulai sekarang, jika $\Gamma_0$ adalah himpunan berhingga !!{sentence},
kita akan menggunakan $\Gamma_0 \Sequent !A$ untuk menyatakan sebarang
sekuen yang antesedennya merupakan barisan !!{sentence} di~$\Gamma_0$,
serta secara tersirat menyertakan kontraksi, pertukaran, dan pelemahan
jika diperlukan.

\begin{defn}[Konsistensi]
Suatu himpunan kalimat~$\Gamma$ \emph{tak konsisten} jika dan hanya jika
terdapat himpunan bagian berhingga~$\Gamma_0 \subseteq \Gamma$ sedemikian
sehingga $\Log{LK}$ !!{derive} $\Gamma_0 \Sequent \quad$. Jika $\Gamma$
tidak tak konsisten, yakni jika untuk setiap $\Gamma_0 \subseteq \Gamma$
yang berhingga, $\Log{LK}$ tidak !!{derive} $\Gamma_0 \Sequent \quad$,
kita menyebutnya \emph{konsisten}.
\end{defn}

\begin{prop}[Refleksivitas]
\ollabel{prop:reflexivity}
Jika $!A \in \Gamma$, maka $\Gamma \Proves !A$.
\end{prop}

\begin{proof}
Sekuen awal $!A \Sequent !A$ !!{derivable}, dan $\{!A\}
\subseteq \Gamma$.
\end{proof}

\begin{prop}[Monotonisitas]
\ollabel{prop:monotonicity}
Jika $\Gamma \subseteq \Delta$ dan $\Gamma \Proves !A$, maka
$\Delta \Proves !A$.
\end{prop}

\begin{proof}
Andaikan $\Gamma \Proves !A$, yakni terdapat $\Gamma_0 \subseteq
\Gamma$ yang berhingga sedemikian sehingga $\Gamma_0 \Sequent !A$
!!{derivable}. Karena $\Gamma \subseteq \Delta$, maka $\Gamma_0$ juga
merupakan himpunan bagian berhingga dari~$\Delta$. Dengan demikian,
!!{derivation} $\Gamma_0 \Sequent !A$ itu juga menunjukkan
$\Delta \Proves !A$.
\end{proof}

\begin{prop}[Transitivitas]
\ollabel{prop:transitivity}
Jika $\Gamma \Proves !A$ dan $\{!A\} \cup \Delta \Proves !B$, maka
$\Gamma \cup \Delta \Proves !B$.
\end{prop}

\begin{proof}
Jika $\Gamma \Proves !A$, terdapat $\Gamma_0 \subseteq \Gamma$ yang
berhingga dan !!a{derivation}~$\pi_0$ untuk $\Gamma_0 \Sequent !A$.
Jika $\{!A\} \cup \Delta \Proves !B$, maka untuk suatu himpunan bagian
berhingga $\Delta_0 \subseteq \Delta$, terdapat !!a{derivation}~$\pi_1$
untuk $!A, \Delta_0 \Sequent !B$. Perhatikan !!{derivation} berikut:
\begin{prooftree}
  \AxiomC{}
  \RightLabel{$\pi_0$}
  \Deduce$\Gamma_0 \fCenter !A$
  \AxiomC{}
  \RightLabel{$\pi_1$}
  \Deduce$!A, \Delta_0 \fCenter !B$
  \RightLabel{\Cut}
  \BinaryInf$\Gamma_0, \Delta_0 \fCenter !B$
\end{prooftree}
Karena $\Gamma_0 \cup \Delta_0 \subseteq \Gamma \cup \Delta$, hal ini
menunjukkan $\Gamma \cup \Delta \Proves !B$.
\end{proof}

Perhatikan bahwa ini khususnya berarti: jika $\Gamma \Proves !A$ dan
$!A \Proves !B$, maka $\Gamma \Proves !B$. Juga diperoleh bahwa jika
$!A_1, \dots, !A_n \Proves !B$ dan $\Gamma \Proves !A_i$ untuk setiap
$i$, maka $\Gamma \Proves !B$.

\begin{prop}
\ollabel{prop:incons}
$\Gamma$ tak konsisten jika dan hanya jika $\Gamma \Proves {!A}$ untuk
setiap kalimat~$!A$.
\end{prop}

\begin{proof}
Latihan.
\end{proof}

\begin{prob}
Buktikan \olref{prop:incons}.
\end{prob}

\begin{prop}[Kekompakan]
  \ollabel{prop:proves-compact}
  \begin{enumerate}
  \item Jika $\Gamma \Proves !A$, maka terdapat himpunan bagian berhingga
    $\Gamma_0 \subseteq \Gamma$ sedemikian sehingga
    $\Gamma_0 \Proves !A$.
  \item Jika setiap himpunan bagian berhingga dari~$\Gamma$ konsisten, maka
    $\Gamma$ konsisten.
  \end{enumerate}
\end{prop}

\begin{proof}
  \begin{enumerate}
    \item Jika $\Gamma \Proves !A$, maka terdapat himpunan bagian berhingga
      $\Gamma_0 \subseteq \Gamma$ sedemikian sehingga sekuen
      $\Gamma_0 \Sequent !A$ mempunyai !!a{derivation}. Akibatnya,
      $\Gamma_0 \Proves !A$.
    \item Jika $\Gamma$ tak konsisten, terdapat himpunan bagian berhingga
      $\Gamma_0 \subseteq \Gamma$ sedemikian sehingga $\Log{LK}$
      !!{derive} $\Gamma_0 \Sequent \quad$. Namun, dengan demikian
      $\Gamma_0$ merupakan himpunan bagian berhingga dari~$\Gamma$ yang tak
      konsisten.
  \end{enumerate}
\end{proof}

\end{document}
